% TODO: 目前内容不足以凑满两页（正反 A4 纸），所以有些内容注释掉了，之后完善

\documentclass{article}
\setlength\parindent{0pt}
\pagenumbering{gobble}

\usepackage{geometry}
\geometry{papersize={210mm,297mm}}
\geometry{left=9mm,right=9mm,top=9mm,bottom=9mm}

\usepackage{fontspec}
\setmainfont{DejaVu Sans}[Scale=0.9]

\usepackage{xeCJK}
\CJKfamily{zhsong}

\usepackage[svgnames]{xcolor}
\usepackage{longfbox}
\usepackage{epstopdf}
\usepackage{float}
\usepackage{xpatch}
\usepackage{tipa}
\usepackage[export]{adjustbox}
\usepackage{enumitem}
\usepackage{mdframed}

\usepackage{multirow}
\usepackage{tabularx}
\renewcommand\tabularxcolumn[1]{m{#1}} % for vertical centering text in X column

\usepackage{multicol}
\usepackage[most]{tcolorbox}
\usepackage{amssymb}
\setlength{\columnsep}{3mm}

\usepackage[explicit,compact]{titlesec}
\titleformat{\section}{\normalfont\bfseries}{}{0pt}{【#1】}

% 啊！万能的 StackExchange！
% https://tex.stackexchange.com/questions/475466/latex-three-column-layout-merging-two-of-them-at-the-begining
%
\newlength{\abstractwidth}
\newlength{\columnshrink}
\newsavebox{\twocolinsert}
%
\makeatletter
\newlength{\resized@col}
\newcounter{column@count}
%
\xpatchcmd{\multi@column@out}{
	\process@cols\mult@gfirstbox{%
		\setbox\count@
		\vsplit\@cclv to\dimen@
		\set@keptmarks
		\setbox\count@
		\vbox to\dimen@
		{\unvbox\count@ \ifshr@nking\vfilmaxdepth\fi}%
	}%
}{
	\process@cols\mult@gfirstbox{%
		\global\advance\c@column@count\@ne
		\resized@col\dimen@%
		\ifnum\c@column@count=\tw@
				\advance\resized@col-\columnshrink
		\fi%
		\setbox\count@
		\vsplit\@cclv to\resized@col
		\set@keptmarks
		\setbox\count@
		\vbox to\dimen@{
			\ifnum
				\c@column@count=\tw@ \vspace*{\columnshrink}
			\fi
			\unvbox\count@
			\ifshr@nking\vfilmaxdepth\fi
		}%
	}%
}{\typeout{Success}}{\typeout{Failure}}
\makeatother

% for designing header
\newsavebox\mysavebox
\newenvironment{imgminipage}[2][]{%
   \def\imgcmd{\includegraphics[width=\wd\mysavebox, height=\dimexpr\ht\mysavebox+\dp\mysavebox\relax, #1]{#2}}%
   \begin{lrbox}{\mysavebox}%
   \begin{minipage}%
}{%
   \end{minipage}
   \end{lrbox}%
   \sbox\mysavebox{\setlength{\fboxrule}{0pt}\fbox{\usebox\mysavebox}}%
   \mbox{\rlap{\raisebox{-\dp\mysavebox}{\imgcmd}}\usebox\mysavebox}%
}

\renewcommand{\labelitemi}{$\blacktriangleright$}

\tcbset{
    frame code={}
    center title,
    left=0pt,
    right=0pt,
    top=6pt,
    bottom=0pt,
    colback=gray!40,
    colframe=white,
    enlarge left by=0mm,
    boxsep=0pt,
    arc=0pt,outer arc=0pt,
}

\begin{document}
\begin{multicols*}{3}

	\setlength{\abstractwidth}{2\linewidth}
	\addtolength{\abstractwidth}{\columnsep}
	\savebox{\twocolinsert}{
	\begin{minipage}{\abstractwidth}
		\noindent 核准日期：2025年05月03日
		\newline 修改日期：2025年05月03日
		\newline

		 \begin{mdframed}[leftline=false, rightline=false, innertopmargin=0pt, innerbottommargin=0pt, innerrightmargin=0pt, innerleftmargin=2em]
		 	\includegraphics[width=0.15\abstractwidth, valign=m]{assets/windows.png}
		 	\hfill
		 	\begin{imgminipage}{assets/header-background.eps}[t]{0.7\abstractwidth}
		 		\Large \textbf{盒装安装媒介说明书}

		 		\normalsize 请仔细阅读说明书并在管理员指导下使用
		 	\end{imgminipage}
		 \end{mdframed}

		% \includegraphics[width=\abstractwidth]{assets/header.eps}


		\begin{mdframed}[hidealllines=true, innerbottommargin=.5em, innertopmargin=0pt]
			\sffamily

			{\centering 警示语 \par}

			无论是否与其它操作系统合用，安装 Microsoft Windows均存在丢失磁盘上所有内容的风险。（参见【不良反应】）

            本说明中的信息和意见无意构成法律建议，法律建议可通过咨询律师来获得。

			一些硬件制造商未能提供基于PnP(Plug-n-Play, 即插即用)协议的设备驱动程序模型，或硬件无法归类于WDF(Windows Driver Framework, 视窗驱动模型)提供的任何一种设备抽象模型。这些硬件无法在Microsoft Windows中正常工作。
		\end{mdframed}
	\end{minipage}}
	\setlength{\columnshrink}{\ht\twocolinsert}
	\addtolength{\columnshrink}{\dp\twocolinsert}
	\noindent\usebox{\twocolinsert}


	\begin{tcolorbox}
	\section*{发行版名称}
	\end{tcolorbox}
	\begin{tabularx}{\linewidth}{@{}ll@{}}
		通用名称： & 视窗操作系统 \\
		正式名称： & Microsoft Windows \\
		拉丁学名： & Fenestrae Molles-Parvae \\
	\end{tabularx}

	\medskip


	\begin{tcolorbox}
	\section*{内容}
	\end{tcolorbox}

	通用 x86-64 Microsoft Windows 基本系统，采用定期发行+滚动升级(Windows Update)模式，尽全力提供最新的稳定版软件，由 Microsoft Corporation 对其进行打包、开发与维护。

	% 内核版本：6.4.1

	% 没有版本号。

	\medskip


	\begin{tcolorbox}
	\section*{性质}
	\end{tcolorbox}

	本系统为采用 Windows NT 内核的操作系统，安装后可由 UEFI 或 Legacy 引导方式启动。

	\medskip


	\begin{tcolorbox}
	\section*{适应平台}
	\end{tcolorbox}

	\begin{itemize}
		\item 仅支持使用 x86-64 (AMD64) 架构的计算机、服务器和嵌入式设备。
		\item 本包装盒中的安装媒介适用于何平台以实际为准。
	\end{itemize}


	\begin{tcolorbox}
	\section*{规格}
	\end{tcolorbox}

	1 枚 安装媒介

	\medskip

	\begin{tcolorbox}
	\section*{用法}
	\end{tcolorbox}

	使用 USB 设备引导。

	启动方式根据硬件调整，一般使用 UEFI。

	根据硬件性能和个人需要，调整安装方式：一般而言，使用 WinPE (Windows Preinstallation  Environment, 视窗预安装环境) 安装。

	安装好基本系统并设置完成软件包后，即可酌情选择可选功能。

	\medskip

	\begin{tcolorbox}
	\section*{不良反应}
	\end{tcolorbox}

	Microsoft Windows使用“微软更新编录”网站提供可以通过公司网络分发的更新列表。您可以将其用作查找Microsoft软件更新、驱动程序和修补程序的一站式位置。每个下载项目都将被分配独立且唯一的KB（Knowledge Base，知识库）编号。

	可以在 https://catalog.update.microsoft.com处得到详细信息。

	\medskip


	\begin{tcolorbox}
	\section*{注意事项}
	\end{tcolorbox}
	\begin{itemize}[leftmargin=*]

		\item 满足最低的硬件要求

		下表是内存和硬盘的需求。

		{\small\begin{tabularx}{\linewidth}{|X|X|X|X|}
			\hline
			类别 & RAM\newline (最低) & RAM\newline (推荐) & 硬盘 \\
			\hline
			完整安装 & 4GB & 8GB & 64GB \\
			\hline
		\end{tabularx}}

		基于您的需求，也许可以使用低于上表所列的配置完成系统安装。但是多数用户在无视这些建议的情况下会安装失败。

		\item 需要固件的设备

		除了需要设备驱动程序，有些硬件还要在使用之前在固件设置中配置工作模式。

	\end{itemize}


	\begin{tcolorbox}
	\section*{禁忌}
	\end{tcolorbox}

	在操作过程中出现或即将出现下列任何一种情况，请立即停止操作，并准备好系统恢复。

	\begin{itemize}[leftmargin=*]
		\setlength{\itemsep}{0pt}
		\setlength{\parskip}{0pt}
		\setlength{\parsep}{0pt}

		\item 以管理员权限在C:\\Windows目录下执行递归删除
		\item 未经确认即删除所有杀毒软件
		\item 未确认设备盘符即执行格式化
		\item 未确认操作即使用命令重定向
		\item 未经确认即执行来自网络的脚本
		\item 长期在散热不良的设备上高负载使用
	\end{itemize}

    % 注释部分就不管了

	% \begin{tcolorbox}
	% \section*{无障碍安装}
	% \end{tcolorbox}

	% Debian GNU/Linux 安装介质不用于视力或运动障碍人士。

	% \medskip

	% \begin{tcolorbox}
	% \section*{新手安装}
	% \end{tcolorbox}

	% Debian GNU/Linux 应谨慎用于新手安装，需要在管理员指导下进行安装，并且需要进行密切的系统监测，一旦出现系统完整度的恶化，应考虑停止使用 Debian GNU/Linux。

	% \medskip


	% \begin{tcolorbox}
	% \section*{版本迭代}
	% \end{tcolorbox}


	\begin{tcolorbox}
	\section*{系统相互作用}
	\end{tcolorbox}
	\begin{itemize}[leftmargin=*]
		\setlength{\parindent}{0pt}

		\item 与 Linux 的相互作用

		当您有双引导时，若另一个操作系统与 Linux 访问相同的文件系统，有可能会导致问题和数据丢失。在这种情况下，文件系统的真实状态可能与 Windows 认为在“启动”之后的情况不同，并且可能在进一步写入文件系统时导致文件系统损坏。因此，在双引导设置中，为了避免文件系统损坏，有必要在 Windows 中禁用“快速启动”功能。

		在罕见情况中已观察到，在使用 Windows 进行系统更新时，可能会出现重新启动后引导被破坏从而导致 Linux 无法启动的情况。同时，若在安装过程中将引导信息写入与 Windows 所在的物理磁盘 MBR 内，将导致后者无法正常启动。

		\item 与其他 Windows 版本的相互作用

		与早于 Windows Vista 的版本共存时，可能导致早期版本的 Windows 无法正常启动。

	\end{itemize}


	\begin{tcolorbox}
	\section*{贮藏}
	\end{tcolorbox}

	-40℃\textasciitilde +70℃

	妥善贮藏所有安装媒介，勿使非安装人员触及。

	\medskip


	\begin{tcolorbox}
	\section*{包装}
	\end{tcolorbox}

	刻录有 Microsoft Windows 安装映像的、兼容 USB 3.0 协议的8GiB闪存驱动器。

	1 枚/盒

	\medskip


	\begin{tcolorbox}
	\section*{有效期}
	\end{tcolorbox}

	尚不明确。

	\medskip


	\begin{tcolorbox}
	\section*{执行标准}
	\end{tcolorbox}
	\begin{tabularx}{\linewidth}{@{}ll@{}}
		\multirow{4}{*}{}{开源许可证：} & GNU GPL（主要）\\
		~ & GNU LGPL \\
		~ & BSD \\
		~ & 以及其他授权条款 \\
	\end{tabularx}

	\medskip


	\begin{tcolorbox}
	\section*{批准文号}
	\end{tcolorbox}

	说明书使用 CC-BY-SA 3.0 协议授权。

	\medskip


% 	\begin{tcolorbox}
% 	\section*{生产单位}
% 	\end{tcolorbox}
%
% 	Debian 计划
%
% 	\medskip


	\begin{tcolorbox}
	\section*{说明书}
	\end{tcolorbox}
	\begin{tabularx}{\linewidth}{@{}ll@{}}
		\multirow{2}{*}{}{编审：} & @YukariChiba\\
		~ & @moesoha \\
        ~ & @Tibrella \\
		图形: & @YJBeetle\\
        ~ & @Isoheptane \\
		GitHub: & isoheptane/arch-media-box\\
	\end{tabularx}

	\medskip


	\vfill
	\begin{flushright}
		Windows \linebreak 11
		\linebreak
		\newline
		\begin{minipage}{0,5\textwidth}
			\centering
			$\vcenter{\hbox{\includegraphics[height=10mm]{assets/microsoft.png}}}$
		\end{minipage}
	\end{flushright}

\end{multicols*}
\end{document}


